\chapter{\texorpdfstring{$\pp$}{p}-Adic Quantum Observable Universe}
\label{chap:p-adic}

% ── Categorical reformulation ───────────────────────────────────────────────

\section{Key Results in Categorical Language}
\label{sec:categorical-reformulation}

\begin{itemize}
  \item Rephrase the QOU definitions
        (Chapter~\ref{chap:qou}) in the language of monoidal categories
        and functors.
  \item Key theorems to establish in the categorical setting:
    \begin{itemize}
      \item Lifting (Proposition~\ref{prop:jet-prolongation}) as a
            categorical lift/extension
            (Remark~\ref{rem:prolongation-categorical}).
      \item Descent (Subsection~\ref{ssec:descent}) via fibred categories.
      \item Obstructions described in terms of cohomology of the category
            (sheaf cohomology, derived functors).
      \item TODO: show that the Atiyah exact sequence splits if and only
            if the bundle admits a linear connection.
    \end{itemize}
  \item Role of $\mathrm{SU}(2)$ and $\mathbb{CP}$ from the symmetry
        perspective; motivate categorical analogies
        (e.g.\ Tannaka--Krein duality).
\end{itemize}

% ── Algebraic structures ────────────────────────────────────────────────────

\section{Universal Enveloping Algebras and Quantum Groups}
\label{sec:enveloping-algebras}

\begin{itemize}
  \item Lusztig integral form of
        $\mathcal{U}_q(\mathfrak{su}(2))$ and its triangular decomposition.
  \item Poincar\'{e}--Birkhoff--Witt (PBW) basis and its
        $q$-deformed analogue.
  \item Universal enveloping algebras and their role in lifting
        (cf.\ \href{https://arxiv.org/abs/math/0409565}{arXiv:math/0409565}).
  \item Heisenberg uncertainty as a $q$-commutator:
        $[\hat{X},\hat{P}]_q = i\hbar_q$.
\end{itemize}

% ── Number fields and universality ──────────────────────────────────────────

\section{\texorpdfstring{$\pp$}{p}-Adic Number Fields and Universality}
\label{sec:p-adic-fields}

\begin{itemize}
  \item Describe $\mathbb{Q}$, $\mathbb{C}_\pp$, $\mathbb{Q}_\pp$, and
        $\mathbb{F}_\pp$ in terms of their category-theoretic universality
        relationships to each other.
  \item Ostrowski's theorem as a categorical coproduct characterisation.
  \item Connections between $\pp$-adic completions and inverse limits in
        the category of topological rings.
\end{itemize}

% ── C_p-QUE ─────────────────────────────────────────────────────────────────

\section{\texorpdfstring{$\mathbb{C}_\pp$}{C\_p} Quantum Universe}
\label{sec:cp-que}

\begin{example}[$\mathbb{Q}_\pp$ realisation]
\label{ex:p-adic-realisation}
Let $k = \mathbb{Q}_\pp$ and let $\mathcal{C}$ be the category of
$F$-isocrystals over~$\mathbb{Q}_\pp$ (modules over the Dieudonn\'e ring
$\mathbb{Q}_\pp\langle F, V \rangle / (FV - \pp)$ that are finite-dimensional
over~$\mathbb{Q}_\pp$).
This is a $\mathbb{Q}_\pp$-linear rigid monoidal category.
The quantum group $\mathcal{O}_q(\mathrm{SU}(2))$ base-changed
to~$\mathbb{Q}_\pp$ provides a Frobenius--Hopf object~$A$, where the
Frobenius form~$\lambda$ arises from the $\pp$-adic Haar measure.
The Frobenius endomorphism~$\varphi$ on $\mathcal{C}$ plays the
r\^ole of the Levi-form contraction: the substrate parameter is
$q = \pp^{\,-v_\pp(\det\varphi)}$ (a $\pp$-adic analogue of
Definition~\ref{def:substrate-parameter}).  The fibre functor~$\omega$
is the forgetful functor to $\mathrm{Vect}_{\mathbb{Q}_\pp}$.
The observability condition
$\mathrm{Hom}(\mathbf{1}, \mathrm{Ad}(A)) = 0$ amounts to the
irreducibility of the associated $F$-crystal, generalising
Lemma~\ref{lem:irreducible-connection} to the $\pp$-adic setting.

\medskip\noindent\emph{Key structures.}
\begin{itemize}
  \item The classical Frobenius integrability theorem constrains
        $q$-torsion leaves to be Legendrian
        (Remark~\ref{rem:q-torsion-existence}).  Over~$\mathbb{Q}_\pp$,
        this is replaced by compatibility of the connection with the
        Frobenius lift: a $\pp$-adic instanton connection is
        ``integrable'' if and only if the associated $F$-crystal is
        unit-root.
  \item Crystal cohomology (de Rham--Witt complex) plays the r\^ole of
        de Rham cohomology.
  \item The base-field-independent
        Definitions~\ref{def:quantum-universe}
        and~\ref{def:quantum-observable-universe} provide the categorical
        framework via Tannakian reconstruction.
\end{itemize}
\end{example}

% ── Finite fields and representations ───────────────────────────────────────

\section{Quotient to Finite Fields}
\label{sec:finite-fields}

\begin{example}[$\mathbb{F}_\pp$ realisation]
\label{ex:finite-field-realisation}
Let $k = \mathbb{F}_\pp$ and let $\mathcal{C} =
\mathrm{Rep}_{\mathbb{F}_\pp}(\mathrm{SU}_q(2))$, the category of
finite-dimensional representations of the quantum group
$\mathrm{SU}_q(2)$ over~$\mathbb{F}_\pp$, where $q \in
\mathbb{F}_\pp^{\times}$.  More generally, when $q$ is a primitive
$\ell$-th root of unity in $\overline{\mathbb{F}}_\pp$ (with $\ell$
coprime to~$\pp$), $\mathcal{C}$ extends to a fusion category over the
extension field $\mathbb{F}_\pp(q)$ with finitely many simple objects.
The Frobenius--Hopf object~$A$ is the regular representation, and
the fibre functor~$\omega$ is the forgetful functor to
$\mathrm{Vect}_{\mathbb{F}_\pp}$.  The non-degeneracy of the
Frobenius form~$\lambda$ follows from the semisimplicity of the fusion
category, and the observability condition
$\mathrm{Hom}(\mathbf{1}, \mathrm{Ad}(A)) = 0$ holds whenever $q \neq 1$
(since the trivial representation does not appear as a summand of the
adjoint representation for generic~$q$).
\end{example}

\begin{itemize}
  \item Quotienting from $\mathbb{Q}_\pp$ to work over $\mathbb{F}_\pp$:
        reduction modulo~$\pp$.
  \item Finite group representations in the QOU context.
  \item Betti numbers over finite fields and comparison with characteristic-zero
        values.
\end{itemize}

% ── Jones polynomials ───────────────────────────────────────────────────────

\section{Jones Polynomials}
\label{sec:jones-polynomials}

% TODO: track the role of Jones polynomials in the $\pp$-adic story and their
% relationship to the braid-group representations of
% Chapter~\ref{chap:path-integrals-braiding}.
